\chapter{Modèles espace-état et filtre de Kalman}\label{ch:qm:state-space-models-and-the-kalman-filter}

Un desk couvre le brut Brent avec le WTI et fixe le ratio de couverture par une régression glissante sur 60 jours des variations quotidiennes du cours du Brent sur celles du WTI. Le ratio
n'a pas été stable~: ajusté année par année, il est passé de 0.41 en 2014 à 0.94 en 2021 et 1.01 en 2026. Lorsque la relation se modifie, la régression glissante
met un mois à couvrir la moitié du mouvement. Un \omterm{def:qm:state-space-models-and-the-kalman-filter:kf}{filtre de Kalman} qui traite le ratio comme un état caché dérivant d'un jour à l'autre promet de
suivre ces déplacements en quelques jours. Ajusté par maximum de vraisemblance sur seize années de cours, le filtre livre une réponse décevante~: le ratio dérive si lentement
par rapport au bruit quotidien que le filtre optimal est lui-même une moyenne sur un mois, et il couvre 0.7~\% mieux que la fenêtre. Un filtre plus rapide
est possible, à un prix que la vraisemblance refuse de payer. Ce chapitre porte sur les \omterm{def:qm:state-space-models-and-the-kalman-filter:ssm}{modèles espace-état} et le filtre qui estime leurs
états cachés~: le modèle linéaire gaussien et les récursions de Kalman, la vraisemblance qu'elles livrent, le lisseur et l'\omterm{def:qm:state-space-models-and-the-kalman-filter:em}{algorithme espérance--maximisation} (EM), les \omterm{def:qm:state-space-models-and-the-kalman-filter:pf}{filtres particulaires} pour
les modèles qui ne sont ni linéaires ni gaussiens, et le ratio de couverture variable dans le temps qui les teste.

\section{Modèles espace-état linéaires gaussiens}

\begin{definition}[Modèle espace-état, modèle à niveau local]\label{def:qm:state-space-models-and-the-kalman-filter:ssm}
Un \emph{modèle espace-état}\index{modèle espace-état} linéaire gaussien comporte une équation d'observation $y_t = Z_t\alpha_t + \varepsilon_t$, $\varepsilon_t \sim \mathcal N(0,
H)$, et une équation d'état $\alpha_{t+1} = T\alpha_t + \eta_t$, $\eta_t \sim \mathcal N(0, Q)$, avec des bruits indépendants et $\alpha_1 \sim \mathcal N(a_1, P_1)$. Le
\emph{modèle à niveau local}\index{modèle à niveau local} est le cas scalaire $Z = T = 1$~: une série observée est égale à un niveau suivant une marche aléatoire, plus un bruit.
\end{definition}

La forme couvre beaucoup de cas~: les modèles ARMA (chapitre~17) dont l'état est formé des valeurs retardées, une régression dont les coefficients dérivent ($Z_t = x_t^\top$,
$T = I$), une structure par terme dont les facteurs évoluent dans le temps, une fair value observée à travers des cotations bruitées. Pour la couverture, $y_t$ est la variation quotidienne du cours du Brent,
$Z_t$ est celle du WTI, et l'état est le ratio de couverture $\beta_t$, une marche aléatoire~: $\Delta b_t = \beta_t\Delta w_t + \varepsilon_t$, $\beta_{t+1} = \beta_t + \eta_t$.

\begin{definition}[Rapport signal sur bruit]\label{def:qm:state-space-models-and-the-kalman-filter:snr}
Dans un \omterm{def:qm:state-space-models-and-the-kalman-filter:ssm}{modèle à niveau local}, ou un modèle à coefficient en marche aléatoire, le \emph{rapport signal sur bruit}\index{rapport signal sur bruit} vaut $q = Q/H$, la variance du
mouvement quotidien de l'état, divisée par la variance du bruit d'observation.
\end{definition}

\section{Le filtre de Kalman}

\begin{definition}[Filtre de Kalman, innovation, gain de Kalman]\label{def:qm:state-space-models-and-the-kalman-filter:kf}
Le \emph{filtre de Kalman}\index{filtre de Kalman} (Kalman, 1960) calcule de manière récursive la loi conditionnelle $\alpha_t \mid y_1, \dots, y_{t-1} \sim \mathcal N(a_t, P_t)$.
Connaissant $(a_t, P_t)$, l'\emph{innovation}\index{innovation} vaut $v_t = y_t - Z_ta_t$, de variance $F_t = Z_tP_tZ_t^\top + H$~; le \emph{gain de Kalman}\index{gain de Kalman}
$K_t = P_tZ_t^\top F_t^{-1}$ met à jour l'état, $a_{t|t} = a_t + K_tv_t$ et $P_{t|t} = P_t - K_tZ_tP_t$, et l'étape de prédiction donne $a_{t+1} = Ta_{t|t}$ et
$P_{t+1} = TP_{t|t}T^\top + Q$.
\end{definition}

\begin{proposition}[Le filtre de Kalman est un conditionnement gaussien]\label{prop:qm:state-space-models-and-the-kalman-filter:kf}
Si $\alpha_t \mid \mathcal F_{t-1} \sim \mathcal N(a_t, P_t)$, alors $\alpha_t \mid \mathcal F_t \sim \mathcal N(a_{t|t}, P_{t|t})$ avec la mise à jour ci-dessus, et
$\alpha_{t+1} \mid \mathcal F_t \sim \mathcal N(a_{t+1}, P_{t+1})$.
\end{proposition}

\begin{proof}
Connaissant $\mathcal F_{t-1}$, $(\alpha_t, y_t)$ est conjointement gaussien, de moyennes $(a_t, Z_ta_t)$, de variances $P_t$ et $F_t$, et de covariance $P_tZ_t^\top$. La
loi conditionnelle d'un vecteur gaussien sachant l'une de ses composantes a pour moyenne $a_t + P_tZ_t^\top F_t^{-1}(y_t - Z_ta_t)$ et pour variance $P_t - P_tZ_t^\top
F_t^{-1}Z_tP_t$. L'étape de prédiction applique l'équation d'état linéaire, avec un bruit indépendant.
\end{proof}

C'est la mise à jour bayésienne du chapitre~14 exécutée une fois par jour, la \omterm{def:qm:bayesian-methods:prior}{loi a posteriori} de la veille servant de \omterm{def:qm:bayesian-methods:prior}{loi a priori} du jour. Une observation manquante se traite en sautant la
mise à jour~: l'étape de prédiction seule élargit $P$. C'est ainsi que le chapitre traite le 20 avril 2020, jour où le prix spot du WTI publié par l'EIA a clôturé à $-36.98$ dollars~:
les variations du WTI de ce jour et du lendemain ($-55.29$ et $+45.89$ dollars, contre $-2.39$ et $-8.24$ pour le Brent) ne relèvent pas d'une relation de couverture mais d'une crise de stockage
(les stocks commerciaux à Cushing, Oklahoma, avaient augmenté de 27 millions de barils depuis mi-mars pour atteindre 83~\% de leur capacité opérationnelle, et le
contrat à terme sur le WTI s'est négocié sous zéro pour la première fois depuis le début de sa cotation en 1983, selon l'EIA), et les deux journées sont marquées comme manquantes.

\begin{proposition}[Gain en régime permanent du modèle à niveau local]\label{prop:qm:state-space-models-and-the-kalman-filter:gain}
Pour le \omterm{def:qm:state-space-models-and-the-kalman-filter:ssm}{modèle à niveau local} de \omterm{def:qm:state-space-models-and-the-kalman-filter:snr}{rapport signal sur bruit} $q$, la variance prédite converge vers $P = pH$ avec $p = \frac12(q + \sqrt{q^2 + 4q})$, et le gain vers
$K = p/(1 + p)$. Le niveau filtré est alors une moyenne exponentiellement pondérée des observations passées, de poids $K$ sur la plus récente.
\end{proposition}

\begin{proof}
En régime permanent $P = P - P^2/(P + H) + Q$, soit $P^2 = Q(P + H)$~; en divisant par $H^2$, $p^2 - qp - q = 0$. La récursion $a_{t+1} = a_t + K(y_t - a_t) =
(1 - K)a_t + Ky_t$ est une moyenne exponentielle.
\end{proof}

Pour $q$ petit, le gain vaut environ $\sqrt q$, de sorte que la mémoire du filtre est d'environ $1/\sqrt q$ jours et sa \omterm{def:qm:stochastic-differential-equations:ou}{demi-vie} $\ln2/\sqrt q$
(\cref{fig:qm:state-space-models-and-the-kalman-filter:gain}). Le \omterm{def:qm:state-space-models-and-the-kalman-filter:snr}{rapport signal sur bruit}, et non l'analyste, détermine la vitesse de réaction du filtre. Pour
un coefficient de régression, il en va de même avec $q$ multiplié par le carré moyen du régresseur.

\begin{omfigure}
\begin{tikzpicture}
\begin{semilogxaxis}[width=10cm,height=5cm,xlabel={rapport signal sur bruit $q = Q/H$},ylabel={gain en régime permanent $K$},
  tick label style={font=\small},label style={font=\small},xmin=1e-4,xmax=100,ymin=0,ymax=1,grid=major,grid style={gray!25},table/col sep=comma]
\addplot[omDef,very thick] table[x=q,y=gain]{figdata/methods/19-state-space-models-and-the-kalman-filter/gain.csv};
\addplot[black,dashed,thin,domain=1e-4:0.3,samples=40] {sqrt(x)};
\end{semilogxaxis}
\end{tikzpicture}

\omcaption{\omterm{def:qm:state-space-models-and-the-kalman-filter:kf}{Gain de Kalman} en régime permanent du \omterm{def:qm:state-space-models-and-the-kalman-filter:ssm}{modèle à niveau local} en fonction du \omterm{def:qm:state-space-models-and-the-kalman-filter:snr}{rapport signal sur bruit}, $K = p/(1 + p)$ avec $p = (q + \sqrt{q^2 + 4q})/2$, et son
approximation pour $q$ petit $\sqrt q$ (pointillés). Pour $q = 0.01$, le filtre accorde 9.5~\% de poids à chaque nouvelle observation~; pour $q = 1$, 62~\%. Données~: la
formule du chapitre.}\label{fig:qm:state-space-models-and-the-kalman-filter:gain}
\end{omfigure}

\begin{definition}[Décomposition par erreurs de prédiction]\label{def:qm:state-space-models-and-the-kalman-filter:ped}
La \emph{décomposition par erreurs de prédiction}\index{décomposition par erreurs de prédiction} écrit la vraisemblance des observations comme produit des densités prédictives à un pas
densités prédictives, $\ln L = -\frac12\sum_t\bigl(\ln(2\pi F_t) + v_t^2/F_t\bigr)$, avec les \omterm{def:qm:state-space-models-and-the-kalman-filter:kf}{innovations} et leurs variances issues du filtre.
\end{definition}

Le filtre fournit donc la vraisemblance de $(H, Q)$ au prix d'un seul passage, et le maximum de vraisemblance ajuste les variances des bruits. Pour la couverture, un raffinement supplémentaire
compte~: les variations quotidiennes des cours du pétrole sont loin d'être homoscédastiques (l'écart-type de celles du WTI valait 0.78 dollars en 2017 et 3.62 en 2026), et une
constante $H$ laisse les années volatiles dicter l'ajustement. Le chapitre pose $\Var(\varepsilon_t) = H s_t^2$, avec $s_t$ une volatilité exponentiellement pondérée des
variations du WTI connue la veille, ce qui revient à filtrer l'équation normalisée $\Delta b_t/s_t = \beta_t\Delta w_t/s_t + \varepsilon_t/s_t$.

\section{Un ratio de couverture variable dans le temps}

Ajusté sur les 4\,097 variations quotidiennes de janvier 2010 à septembre 2026, le modèle normalisé donne $H = 0.508$, $Q = 0.000238$ et un \omterm{def:qm:state-space-models-and-the-kalman-filter:snr}{rapport signal sur bruit} de
0.00047~; avec un carré moyen du régresseur normalisé de 1.12, la demi-réponse du filtre à un déplacement vaut $\ln2/\sqrt{0.00047 \times 1.12} = 30$ jours, soit exactement les
30 jours qu'il faut à une fenêtre de 60 jours pour couvrir la moitié d'un déplacement. L'incertitude du ratio filtré est elle-même élevée~: son écart-type vaut environ 0.10, si bien que le ratio quotidien
n'est connu qu'à $\pm 0.2$ près (\cref{fig:qm:state-space-models-and-the-kalman-filter:ratio}).

\begin{omfigure}
\begin{tikzpicture}
\begin{axis}[width=12cm,height=5.4cm,xlabel={année},ylabel={ratio de couverture (Brent sur WTI)},
  tick label style={font=\small},label style={font=\small},xmin=2019,xmax=2026.8,ymin=0,ymax=1.6,grid=major,grid style={gray!25},
  xticklabel style={/pgf/number format/set thousands separator={}},
  legend columns=4,legend style={font=\footnotesize,at={(0.5,-0.3)},anchor=north,draw=none,fill=none,
  /tikz/every even column/.append style={column sep=6pt}},table/col sep=comma]
\addplot[name path=hi,draw=none,forget plot] table[x=year,y=hi]{figdata/methods/19-state-space-models-and-the-kalman-filter/ratio.csv};
\addplot[name path=lo,draw=none,forget plot] table[x=year,y=lo]{figdata/methods/19-state-space-models-and-the-kalman-filter/ratio.csv};
\addplot[omDef!15] fill between[of=hi and lo];
\addplot[omDef,thick] table[x=year,y=kalman]{figdata/methods/19-state-space-models-and-the-kalman-filter/ratio.csv};
\addplot[black,very thick] table[x=year,y=smooth]{figdata/methods/19-state-space-models-and-the-kalman-filter/ratio.csv};
\addplot[omThm,thin] table[x=year,y=roll]{figdata/methods/19-state-space-models-and-the-kalman-filter/ratio.csv};
\legend{$\pm 2$ écarts-types,filtre de Kalman,lisseur de Kalman,régression glissante sur 60 jours}
\end{axis}
\end{tikzpicture}

\omcaption{Ratio de couverture Brent contre WTI sur les variations quotidiennes des prix, 2019--2026~: le ratio filtré par Kalman avec deux écarts-types de son incertitude, le
ratio lissé (qui utilise l'ensemble de l'échantillon), et la régression glissante sur 60 jours. Données~: séries FRED DCOILBRENTEU et DCOILWTICO (US Energy Information
Administration, prix spot).}\label{fig:qm:state-space-models-and-the-kalman-filter:ratio}
\end{omfigure}

Chaque estimation couvre le jour $t$ avec le ratio connu à la clôture du jour $t - 1$. Sur les 4\,035 jours suivant un démarrage de 60 jours, l'erreur quadratique moyenne de couverture, en
dollars carrés par baril et par jour, vaut 3.324 sans couverture, 1.329 avec le ratio des 60 premiers jours maintenu fixe, 1.276 avec la régression sur 60 jours et 1.267 avec
le \omterm{def:qm:state-space-models-and-the-kalman-filter:kf}{filtre de Kalman}~: une réduction de 0.7~\% par rapport à la fenêtre. Sans la mise à l'échelle de la volatilité, le filtre du maximum de vraisemblance est \emph{moins bon} de 1.4~\% que la
fenêtre, car les années volatiles multiplient son $Q$ par dix. Rendre le filtre plus rapide est possible~: multiplier $Q$ par 10 ramène le temps de demi-réponse à 9.6 jours et
augmente l'erreur de couverture de 0.6~\% par rapport à la fenêtre~; la multiplier par 100 donne 3.0 jours et 6.2~\% d'erreur en plus
(la \cref{fig:qm:state-space-models-and-the-kalman-filter:tradeoff}). Dans une simulation où le ratio passe de 0.7 à 1.0 avec les niveaux de bruit estimés, le
filtre du maximum de vraisemblance et la fenêtre ont besoin d'environ cinquante jours pour atteindre 0.85~; celui avec un $Q$ cent fois plus grand n'en demande que huit.

\begin{omfigure}
\begin{tikzpicture}
\begin{semilogxaxis}[width=10cm,height=5cm,xlabel={jours de couverture pour une demi-variation du ratio},ylabel={erreur de couverture relative},
  tick label style={font=\small},label style={font=\small},xmin=2,xmax=150,ymin=0.97,ymax=1.08,grid=major,grid style={gray!25},
  xtick={3,10,30,100},xticklabels={3,10,30,100},
  yticklabel style={/pgf/number format/fixed,/pgf/number format/precision=2},
  legend style={font=\footnotesize,at={(0.97,0.97)},anchor=north east,fill=white,draw=none},table/col sep=comma]
\addplot[omDef,very thick,mark=*,mark size=1.8pt] table[x=halfdays,y=relvar]{figdata/methods/19-state-space-models-and-the-kalman-filter/tradeoff.csv};
\addplot[only marks,mark=square*,mark size=2.5pt,omThm] coordinates {(30,1.0)};
\addplot[only marks,mark=o,mark size=4pt,black,thick] coordinates {(30.26,0.9928)};
\legend{filtre de Kalman ($Q \times 0.1$ à $\times 100$),régression sur 60 jours,maximum de vraisemblance}
\end{semilogxaxis}
\end{tikzpicture}

\omcaption{Vitesse contre précision pour la couverture Brent--WTI, 2010--2026~: chaque point est un \omterm{def:qm:state-space-models-and-the-kalman-filter:kf}{filtre de Kalman} dont le $Q$ estimé est multiplié par 0.1, 0.3, 1, 3, 10, 30 et
100, placé selon son temps de demi-réponse et son erreur quadratique moyenne de couverture relative à la régression sur 60 jours. Le choix du maximum de vraisemblance (entouré) se place à côté de la
fenêtre~; un filtre qui suit les variations à l'échelle de quelques jours le paie en bruit. Données~: séries FRED DCOILBRENTEU et DCOILWTICO (EIA).}\label{fig:qm:state-space-models-and-the-kalman-filter:tradeoff}
\end{omfigure}

\section{Le lisseur et l'algorithme espérance--maximisation}

\begin{definition}[Lisseur de Kalman]\label{def:qm:state-space-models-and-the-kalman-filter:smoother}
Le \emph{lisseur de Kalman}\index{lisseur de Kalman} calcule $\E[\alpha_t \mid y_1, \dots, y_n]$ et sa variance à partir de la sortie du filtre, par une passe rétrograde
(Rauch, Tung et Striebel, 1965)~: $\hat\alpha_t = a_{t|t} + J_t(\hat\alpha_{t+1} - a_{t+1})$ avec $J_t = P_{t|t}T^\top P_{t+1}^{-1}$.
\end{definition}

Le lisseur répond à une question différente de celle du filtre~: non pas ce que le desk pouvait savoir, mais quelle valeur le ratio avait le plus vraisemblablement. C'est l'outil de
recherche sur données passées (la \cref{fig:qm:state-space-models-and-the-kalman-filter:ratio} le montre plus lisse et plus précoce que le filtre) et une source de biais d'anticipation
(look-ahead) lorsqu'il s'infiltre dans un backtest.

\begin{definition}[Algorithme espérance--maximisation]\label{def:qm:state-space-models-and-the-kalman-filter:em}
L'\emph{algorithme espérance--maximisation}\index{algorithme espérance--maximisation} (Dempster, Laird et Rubin, 1977) maximise une vraisemblance avec
variables non observées, en alternant une étape E, l'espérance de la log-vraisemblance des données complètes sachant les données et les paramètres courants, et une étape M qui
la maximise. Pour les \omterm{def:qm:state-space-models-and-the-kalman-filter:ssm}{modèles espace-état}, l'étape E est le lisseur et l'étape M pour $H$ et $Q$ admet une forme fermée (Shumway et Stoffer, 1982).
\end{definition}

Chaque étape EM ne peut faire diminuer la vraisemblance~: l'espérance de la log-vraisemblance des données complètes est une borne inférieure de la log-vraisemblance, à une constante près, qui
touche celle-ci aux paramètres courants, et l'étape M relève la borne. EM est stable et n'exige aucune dérivée, mais il peut être lent. Lancé à partir de $H = 1$ et $Q =
0.01$ sur le modèle de couverture, elle fait monter la log-vraisemblance de $-4\,947$ à $-4\,524$ en dix étapes et à $-4\,502$ en quarante, soit encore 32 points sous le maximum
$-4\,470$ que la maximisation directe atteint (la \cref{fig:qm:state-space-models-and-the-kalman-filter:em})~: la vraisemblance est plate le long de la direction qui échange $Q$
contre $H$.

\begin{omfigure}
\begin{tikzpicture}
\begin{axis}[width=10cm,height=4.8cm,xlabel={itération EM},ylabel={log-vraisemblance},
  tick label style={font=\small},label style={font=\small},xmin=0,xmax=41,ymin=-4600,ymax=-4450,grid=major,grid style={gray!25},
  yticklabel style={/pgf/number format/set thousands separator={\,}},
  legend style={font=\footnotesize,at={(0.97,0.05)},anchor=south east,fill=white,draw=none},table/col sep=comma]
\addplot[omDef,very thick,mark=*,mark size=1pt] table[x=iter,y=loglik]{figdata/methods/19-state-space-models-and-the-kalman-filter/em.csv};
\addplot[black,dashed,thick] table[x=iter,y=mle]{figdata/methods/19-state-space-models-and-the-kalman-filter/em.csv};
\legend{{EM à partir de $H = 1$, $Q = 0.01$},maximum de vraisemblance}
\end{axis}
\end{tikzpicture}

\omcaption{EM sur le modèle de couverture Brent--WTI mis à l'échelle~: la log-vraisemblance croît à chaque étape (depuis $-4\,947$ au départ, hors du graphique) mais s'approche lentement du maximum
trouvé par optimisation directe (en tirets, $-4\,470$). Données~: séries FRED DCOILBRENTEU et DCOILWTICO (EIA).}\label{fig:qm:state-space-models-and-the-kalman-filter:em}
\end{omfigure}

\section{Filtres particulaires}

\begin{definition}[Filtre particulaire, rééchantillonnage d'importance séquentiel]\label{def:qm:state-space-models-and-the-kalman-filter:pf}
Un \emph{filtre particulaire}\index{filtre particulaire} représente la loi de filtrage d'un état par un échantillon pondéré de particules. \emph{Rééchantillonnage
d'importance séquentiel}\index{rééchantillonnage d'importance séquentiel}, le filtre bootstrap de Gordon, Salmond et Smith (1993), fait progresser chaque particule à travers l'équation
d'état, pondère chaque particule par la densité de la nouvelle observation sachant la particule, et rééchantillonne en proportion des poids~; le poids moyen à chaque étape
estime la densité prédictive à un pas, et leur produit la vraisemblance.
\end{definition}

Le \omterm{def:qm:state-space-models-and-the-kalman-filter:pf}{filtre particulaire} n'exige ni linéarité ni bruit gaussien, seulement la capacité de simuler l'état et d'évaluer la densité d'observation~; son prix est
une erreur de Monte Carlo et une dégénérescence, mesurée par la \omterm{def:qm:bayesian-methods:ess}{taille d'échantillon effective} du chapitre~14 calculée à partir des poids, $1/\sum_iw_i^2$. Sur un niveau local linéaire gaussien
de 500 observations, où le \omterm{def:qm:state-space-models-and-the-kalman-filter:kf}{filtre de Kalman} est exact, 2\,000 particules donnent une log-vraisemblance de $-785.06$, contre la valeur exacte de $-784.22$, et des moyennes filtrées
à 0.07 près des valeurs exactes (dont l'écart-type a posteriori vaut 0.45), avec une \omterm{def:qm:bayesian-methods:ess}{taille d'échantillon effective} moyenne de 1\,678.

Le cas utile est le cas non linéaire. Un modèle de volatilité stochastique pour les rendements quotidiens EUR/USD laisse la log-variance suivre un AR(1), $x_t = \mu + \phi(x_{t-1} - \mu)
+ \sigma_\eta u_t$, et observe $r_t = e^{x_t/2}z_t$~: aucun \omterm{def:qm:state-space-models-and-the-kalman-filter:kf}{filtre de Kalman} ne s'applique. Avec 5\,000 particules sur les 1\,000 derniers rendements, une recherche sur grille de la
vraisemblance particulaire (mêmes nombres aléatoires en chaque point) retient $\phi = 0.9$ et $\sigma_\eta = 0.4$ (log-vraisemblance $-568.4$, contre $-572.3$ pour $\phi = 0.95$,
$\sigma_\eta = 0.2$~: une surface plate). La volatilité filtrée suit la volatilité GARCH du chapitre~18, mais réagit plus vite
(la \cref{fig:qm:state-space-models-and-the-kalman-filter:sv})~; la \omterm{def:qm:bayesian-methods:ess}{taille d'échantillon effective} est en moyenne de 4\,449, mais tombe à 23 le 3 avril 2025, lorsque l'euro a progressé de 2.7~\%
après l'annonce de droits de douane américains du 2 avril (épisode d'aversion pour le risque où, fait rare, le dollar s'est déprécié, comme le relève le Financial Stability Review de la BCE de novembre
2025) et presque toutes les particules étaient trop calmes pour l'expliquer.

\begin{omfigure}
\begin{tikzpicture}
\begin{axis}[width=12cm,height=5cm,xlabel={année},ylabel={volatilité quotidienne (\%)},
  tick label style={font=\small},label style={font=\small},xmin=2022.8,xmax=2026.8,ymin=0,ymax=1.5,grid=major,grid style={gray!25},
  xtick={2023,2024,2025,2026},xticklabel style={/pgf/number format/set thousands separator={}},
  legend columns=2,legend style={font=\footnotesize,at={(0.5,-0.3)},anchor=north,draw=none,fill=none,
  /tikz/every even column/.append style={column sep=8pt}},table/col sep=comma]
\addplot[omDef,thick] table[x=year,y=sv]{figdata/methods/19-state-space-models-and-the-kalman-filter/sv.csv};
\addplot[omThm,very thick,dashed] table[x=year,y=garch]{figdata/methods/19-state-space-models-and-the-kalman-filter/sv.csv};
\legend{{volatilité stochastique, filtre particulaire},GARCH-$t$ (chapitre~18)}
\end{axis}
\end{tikzpicture}

\omcaption{Volatilité quotidienne EUR/USD sur les 1\,000 derniers fixings de la BCE jusqu'au 23 septembre 2026~: l'estimation de volatilité stochastique filtrée par particules ($\phi = 0.9$,
$\sigma_\eta = 0.4$, 5\,000 particules) et la volatilité conditionnelle GARCH-$t$. Données~: taux de référence de l'euro de la BCE (source~: statistiques
de la BCE).}\label{fig:qm:state-space-models-and-the-kalman-filter:sv}
\end{omfigure}

\section{Tutoriel~: le ratio de couverture qui a bougé}\label{tut:qm:state-space-models-and-the-kalman-filter:hedge}

\begin{tutorial}
\textbf{Objectif.} Suivre le ratio de couverture Brent--WTI à l'aide d'un \omterm{def:qm:state-space-models-and-the-kalman-filter:kf}{filtre de Kalman} estimé par maximum de vraisemblance, et mesurer ce qu'il apporte face à une régression glissante.
\textbf{État final~:} les \cref{fig:qm:state-space-models-and-the-kalman-filter:ratio,fig:qm:state-space-models-and-the-kalman-filter:tradeoff}, le \omterm{def:qm:state-space-models-and-the-kalman-filter:snr}{rapport signal sur bruit}
de 0.00047, et la réduction de 0.7~\% de l'erreur de couverture.
\begin{tutsteps}
\item \textbf{Le filtre}, avec observations manquantes et vraisemblance fondée sur la \omterm{def:qm:state-space-models-and-the-kalman-filter:ped}{décomposition par erreurs de prédiction}.
\omcode{code/firm/kalman/firm_kalman.py}{32}{58}{Le filtre de Kalman avec la décomposition par erreurs de prédiction.}\label{lst:qm:state-space-models-and-the-kalman-filter:kf}
\item \textbf{Le suiveur en flux continu} que le desk lancerait chaque soir.
\omcode{code/firm/kalman/firm_kalman.py}{174}{188}{Un ratio de couverture variable dans le temps, en flux continu.}\label{lst:qm:state-space-models-and-the-kalman-filter:tracker}
\item \textbf{Exécuter} \texttt{compare()}, \texttt{tradeoff()}, \texttt{em\_path()} et \texttt{sv\_grid()} dans \texttt{qm\_kalman.py}, puis \texttt{fig\_kalman.py}.
\end{tutsteps}
\textbf{Pour aller plus loin.} Laisser le ratio revenir vers sa moyenne ($T = \phi < 1$ avec une moyenne) au lieu de le laisser dériver, et comparer les vraisemblances~; ajouter la variation retardée du WTI comme
second régresseur pour absorber le décalage temporel entre les deux cours au comptant.
\end{tutorial}

\section{Construction~: la boîte à outils de Kalman}\label{bld:qm:state-space-models-and-the-kalman-filter:kalman}

\begin{build}
\bfield{Rôle} Les grandeurs cachées que la firme miniature suit chaque jour (ratios de couverture, fair values observées à travers des prix bruités, expositions aux signaux qui dérivent) sont
filtrées ici, avec leur incertitude.
\bfield{Interface} \texttt{kalman\_filter(y, Z, T, H, Q, a0, P0)}~; \texttt{rts\_smoother(kf, T)}~; \texttt{em(\dots)}~; \texttt{fit\_mle(\dots)}~;
\texttt{local\_level\_gain(q)}~; \texttt{particle\_filter(y, n, init, step, loglik\_obs, rng)}~; \texttt{HedgeTracker(q, h, beta0, p0).update(x, y)}.
\bfield{Règles} Les filtres utilisés pour le trading n'emploient que les données jusqu'à la clôture précédente~; les lisseurs servent à la recherche et n'alimentent jamais un backtest~; les données manquantes sont notées NaN, et non à zéro~;
chaque \omterm{def:qm:state-space-models-and-the-kalman-filter:pf}{filtre particulaire} publie sa \omterm{def:qm:bayesian-methods:ess}{taille d'échantillon effective}.
\bfield{Tests d'acceptation} \texttt{code/firm/kalman/tests/}~: la vraisemblance et les états lissés égalent le conditionnement gaussien direct sur une courte série~; une
observation manquante saute la mise à jour~; la formule du gain en régime permanent~; le maximum de vraisemblance et l'EM (monotone) retrouvent un niveau local simulé~; le \omterm{def:qm:state-space-models-and-the-kalman-filter:pf}{filtre particulaire}
concorde avec le \omterm{def:qm:state-space-models-and-the-kalman-filter:kf}{filtre de Kalman} sur un modèle linéaire~; le tracker en flux continu reproduit le filtre.
\bfield{Pour aller plus loin} Filtrage par racine carrée pour la stabilité numérique~; les filtres étendu et sans parfum~; Metropolis--Hastings marginal particulaire pour les
paramètres de la volatilité stochastique.
\end{build}

\begin{omsources}
\item EIA, \emph{Today in Energy}, sur les cours du pétrole brut en 2020 (WTI négatif le 20 avril 2020)~; BCE, \emph{Financial Stability Review}, novembre 2025, encadré spécial sur l'annonce de droits de douane d'avril 2025.
\item R.~E.~Kalman, ``A new approach to linear filtering and prediction problems'', \emph{Journal of Basic Engineering} 82, 1960.
\item H.~E.~Rauch, F.~Tung and C.~T.~Striebel, ``Maximum likelihood estimates of linear dynamic systems'', \emph{AIAA Journal} 3, 1965.
\item A.~P.~Dempster, N.~M.~Laird et D.~B.~Rubin, ``Maximum likelihood from incomplete data via the EM algorithm'', \emph{Journal of the Royal
Statistical Society B} 39, 1977.
\item R.~H.~Shumway et D.~S.~Stoffer, ``An approach to time series smoothing and forecasting using the EM algorithm'', \emph{Journal of Time Series
Analysis} 3, 1982.
\item N.~J.~Gordon, D.~J.~Salmond et A.~F.~M.~Smith, ``Novel approach to nonlinear/non-Gaussian Bayesian state estimation'', \emph{IEE Proceedings F} 140,
1993.
\item US Energy Information Administration, cours au comptant du WTI et du Brent, via FRED (séries DCOILWTICO et DCOILBRENTEU), consulté le 24 septembre 2026.
\end{omsources}

\section{Exercices}

\begin{exercise}[$\star$]\label{exo:qm:state-space-models-and-the-kalman-filter:1}
Soit un \omterm{def:qm:state-space-models-and-the-kalman-filter:ssm}{modèle à niveau local} avec $H = 1$ et $Q = 0.04$. Calculer le gain en régime permanent et la \omterm{def:qm:stochastic-differential-equations:ou}{demi-vie} de la mémoire du filtre.
\end{exercise}

\begin{exercise}[$\star$]\label{exo:qm:state-space-models-and-the-kalman-filter:2}
Le ratio de couverture filtré de la veille vaut $0.80$, de variance $0.01$~; $Q = 0.0002$, $H = 0.5$. Aujourd'hui, le WTI varie de $+2$ et le Brent de $+1.2$ (unités normalisées). Mettre à jour le ratio.
\end{exercise}

\begin{exercise}[$\star$]\label{exo:qm:state-space-models-and-the-kalman-filter:3}
Expliquer pourquoi une observation manquante augmente la variance de l'état, et de combien dans le \omterm{def:qm:state-space-models-and-the-kalman-filter:ssm}{modèle à niveau local}.
\end{exercise}

\begin{exercise}[$\star\star$]\label{exo:qm:state-space-models-and-the-kalman-filter:4}
Montrer que, pour $q$ petit, le gain du niveau local vaut environ $\sqrt q$, et comparer avec une moyenne mobile exponentielle~;
déterminer le $\lambda$ que le filtre implique pour $q = 0.0004$.
\end{exercise}

\begin{exercise}[$\star\star$]\label{exo:qm:state-space-models-and-the-kalman-filter:5}
Écrire le modèle AR(1) plus bruit, $y_t = x_t + \varepsilon_t$ avec $x_t = \phi x_{t-1} + \eta_t$, sous forme espace-état, et donner son équation
de la variance prédite en régime permanent.
\end{exercise}

\begin{exercise}[$\star\star$]\label{exo:qm:state-space-models-and-the-kalman-filter:6}
Les poids normalisés d'un \omterm{def:qm:state-space-models-and-the-kalman-filter:pf}{filtre particulaire} valent 0.5, 0.3, 0.1, 0.05 et 0.05. Calculer la \omterm{def:qm:bayesian-methods:ess}{taille d'échantillon effective}.
\end{exercise}

\begin{exercise}[$\star\star\star$]\label{exo:qm:state-space-models-and-the-kalman-filter:7}
\emph{Programmation.} Sur les données Brent--WTI, multiplier le $Q$ estimé par 0.1, 1, 10 et 100 et tabuler la variance de l'erreur de couverture par rapport à la régression sur 60 jours et
le temps de demi-réponse.
\end{exercise}

\begin{exercise}[$\star\star\star$]\label{exo:qm:state-space-models-and-the-kalman-filter:8}
\emph{Trouver la faille.} «~Nous avons soumis une couverture à un backtest avec le ratio lissé par Kalman et réduit de 15~\% l'erreur de couverture par rapport à la régression glissante.~»
\end{exercise}

\section{Problème~: le ratio de couverture qui a bougé}

\begin{problem}[{Problème du week-end --- à quelle vitesse doit varier un ratio de couverture~?}]\label{pb:qm:state-space-models-and-the-kalman-filter:1}
Les données sont les cours quotidiens au comptant du WTI et du Brent de l'EIA, obtenus via FRED, du 4 janvier 2010 au 22 septembre 2026 (4\,097 variations quotidiennes). La variation quotidienne du cours du Brent est
couverte avec celle du WTI~; le ratio de couverture du jour $t$ doit être connu à la clôture du jour $t - 1$.

\medskip\noindent\textbf{Partie I --- Les données.}
\begin{enumerate}
\item Que s'est-il passé le 20 avril 2020, et comment le chapitre le traite-t-il~?
\item Comment le ratio de régression annuel a-t-il évolué entre 2014, 2021 et 2026~?
\item Pourquoi normaliser le bruit d'observation par la volatilité du WTI~?
\item Quelles sont les variances de l'erreur non couverte et de l'erreur de couverture à ratio fixe~?
\item Quel résultat obtient la régression sur 60 jours~?
\end{enumerate}

\medskip\noindent\textbf{Partie II --- Le filtre.}
\begin{enumerate}[resume]
\item Quels sont $H$, $Q$ et le \omterm{def:qm:state-space-models-and-the-kalman-filter:snr}{rapport signal sur bruit} donnés par le maximum de vraisemblance~?
\item Quelle est la demi-réponse du filtre à un choc, et comment se compare-t-elle à celle de la fenêtre~?
\item Quelle variance de l'erreur de couverture le filtre atteint-il, et quelle réduction par rapport à la fenêtre~?
\item Que se passe-t-il sans la normalisation par la volatilité~?
\item Avec quelle précision le ratio quotidien est-il connu~?
\end{enumerate}

\medskip\noindent\textbf{Partie III --- Vitesse et bruit.}
\begin{enumerate}[resume]
\item Que donnent les filtres dont $Q$ est multiplié par 10 et par 100~?
\item Dans le saut simulé de 0.7 à 1.0, combien de jours les filtres et la fenêtre mettent-ils pour atteindre 0.85~?
\item Que trouve l'EM à partir d'un mauvais point de départ, et pourquoi est-il lent~?
\item Quels sont les ratios filtré et glissant en fin d'échantillon~?
\item Quand un filtre rapide compense-t-il son bruit~?
\end{enumerate}

\medskip\noindent\textbf{Partie IV --- Jugement.}
\begin{enumerate}[resume]
\item Le desk doit-il remplacer sa régression glissante~?
\item Qu'est-ce qui pourrait améliorer la couverture davantage que le choix du filtre~?
\item Pourquoi le lisseur doit-il rester en dehors du backtest~?
\item Énoncer le \emph{résultat nommé}~: la réduction de la variance de l'erreur de couverture et le \omterm{def:qm:state-space-models-and-the-kalman-filter:snr}{rapport signal sur bruit} estimé.
\item En une phrase~: qu'est-ce qui détermine la vitesse de réaction d'un \omterm{def:qm:state-space-models-and-the-kalman-filter:kf}{filtre de Kalman}~?
\end{enumerate}
\end{problem}

\section{Questions d'entretien}

\begin{interviewq}[$\star$ \iqroles{chercheur, trader}]\label{iq:qm:state-space-models-and-the-kalman-filter:1}
Expliquez le \omterm{def:qm:state-space-models-and-the-kalman-filter:kf}{filtre de Kalman} à un trader en trois phrases.
\end{interviewq}

\begin{interviewq}[$\star\star$ \iqroles{chercheur}]\label{iq:qm:state-space-models-and-the-kalman-filter:2}
Établissez la mise à jour de Kalman pour un état scalaire à partir du conditionnement gaussien.
\end{interviewq}

\begin{interviewq}[$\star\star$ \iqroles{chercheur, trader}]\label{iq:qm:state-space-models-and-the-kalman-filter:3}
Comment choisissez-vous le bruit de processus d'un \omterm{def:qm:state-space-models-and-the-kalman-filter:kf}{filtre de Kalman} qui suit un ratio de couverture~?
\end{interviewq}

\begin{interviewq}[$\star\star$ \iqroles{chercheur, ingénieur ML}]\label{iq:qm:state-space-models-and-the-kalman-filter:4}
Quelle est la différence entre filtrage et lissage, et pourquoi cela compte-t-il dans les backtests~?
\end{interviewq}

\begin{interviewq}[$\star\star$ \iqroles{ingénieur ML, chercheur}]\label{iq:qm:state-space-models-and-the-kalman-filter:5}
Quand utiliseriez-vous un \omterm{def:qm:state-space-models-and-the-kalman-filter:pf}{filtre particulaire} plutôt qu'un \omterm{def:qm:state-space-models-and-the-kalman-filter:kf}{filtre de Kalman}, et que peut-il mal se passer~?
\end{interviewq}

\begin{interviewq}[$\star\star\star$ \iqroles{chercheur}]\label{iq:qm:state-space-models-and-the-kalman-filter:6}
Pourquoi chaque itération de l'EM ne fait-elle pas diminuer la vraisemblance~?
\end{interviewq}
